% Copy each caption into the matching figure environment.
% PDF exports are tightly cropped. Use width=\textwidth in a figure* environment.

% fig1_wrapper_effects.pdf
\caption{Wrapper effects for attributed false assertions. Points show the mean
change in answer margin from the same paragraph without a wrapper. Bars show
95\% paired bootstrap confidence intervals over facts. A positive change means
greater relative preference for the assigned false answer. Each tag is next to
its colon-label form. The two models use different horizontal scales. Within
each model, both datasets use the same scale.}
\label{fig:wrapper-effects}

% fig2_name_gaps.pdf
\caption{Tag-versus-colon differences across names. Each coloured point is one
name's mean answer margin under tags minus its mean margin under a colon label.
Positive values mean greater relative preference for the false answer under
tags. Each panel shows 30 random names and the six tested meaningful names.
Black diamonds show group means. Horizontal intervals show the reported
95\% crossed bootstrap confidence intervals over names and facts for the
random-name means. No group-level interval is shown for the six meaningful
names. Vertical offsets separate overlapping points and have no numeric
meaning. The two models use different horizontal scales.}
\label{fig:name-gaps}

% fig3_instruction_guards.pdf
\caption{Warnings against an embedded instruction inside a
\texttt{<document>} wrapper. Cells show the percentage of outputs matching the
assigned false answer. The no-warning rate is 100\% in all four panels.
Each warning text is placed in the system prompt, above the block in the user
turn, or below the block in the user turn. All panels use the same colour scale.}
\label{fig:instruction-guards}

% appendix_assertion_guards.pdf
\caption{Warnings against an attributed false assertion inside a
\texttt{<document>} wrapper. Cells show the percentage of outputs matching the
assigned false answer. The no-warning row is repeated across columns for
reference. All panels use the same colour scale as
Figure~\ref{fig:instruction-guards}.}
\label{fig:assertion-guards}
